\documentclass[11pt]{article}
\usepackage[a4paper,margin=1in]{geometry}
\usepackage{amsmath,amssymb,amsthm}
\usepackage{booktabs}
\usepackage{array}
\usepackage{hyperref}

\newtheorem{definition}{Definition}
\newtheorem{theorem}{Theorem}
\newtheorem{nogo}{No-Go}

\newcommand{\R}{\mathbb R}
\newcommand{\Hp}{H_{\mathrm{HP}}}
\newcommand{\Hxp}{H_{xp}}
\newcommand{\XiHP}{\Xi_{\mathrm{HP}}}
\newcommand{\XiBC}{\Xi_{\mathrm{BC}}}
\newcommand{\XidB}{\Xi_{\mathrm{dB}}}

\title{Step 182: Hilbert--Polya / Berry--Keating Carrier Pivot}
\author{RH Membrane Adequacy Track}
\date{}

\begin{document}
\maketitle

\section*{Verdict}

\[
\boxed{\mathrm{HP\_verdict}=V_{\mathrm{HP\_meta\_pattern\_confirmed}}.}
\]

Carrier-specific outcome: the standard Berry--Keating operator on the
half-line has the wrong spectral type, while exact modified spectral matching
is the Hilbert--Polya/RH assertion itself.  Thus the HP/BK pivot contributes a
new typed no-go:
\[
\boxed{\text{Hilbert--Polya/Berry--Keating standard \(xp\) bridge failure and modified-operator target-equivalence}.}
\]

\section*{Inherited Record Audit}

The local inherited records contain no accepted Hilbert--Polya or
Berry--Keating carrier package.  A search through the active RH membrane
records, \texttt{anti\_loc/adequacy.tex}, and \texttt{anti\_loc/needles.tex}
does not provide an \(H_{xp}\) carrier, a self-adjoint extension ledger, or a
spectral determinant bridge to \(\xi\).

The closest inherited record is Step 93's Tier-0 triage.  It lists
Connes--Consani--Moscovici finite spectral triples as strong finite evidence,
but explicitly warns that convergence to the completed determinant or zero
ledger remains the true gate.  Step 99 repeats the same promotion rule:
finite spectral windows are support-only unless an exhaustive tail or
completed ledger record is supplied.

\section*{Carrier Declaration}

\begin{definition}[Standard HP/BK half-line carrier]
Choose
\[
\Hp=L^2(\R_+,dx),\qquad
\langle f,g\rangle=\int_0^\infty f(x)\overline{g(x)}\,dx.
\]
The symmetrized Berry--Keating operator is
\[
\Hxp=\frac{xp+px}{2}
     =-i\left(x\frac{d}{dx}+\frac12\right)
\]
on the core \(C_c^\infty(\R_+)\).
\end{definition}

The choice \(L^2(\R_+,dx)\) is deliberate: it is the carrier on which the
operator has exactly the symmetrized form written above.  The logarithmic
unitary
\[
(Uf)(t)=e^{t/2}f(e^t),\qquad
U:L^2(\R_+,dx)\to L^2(\R,dt),
\]
gives
\[
U\Hxp U^{-1}=-i\frac{d}{dt}.
\]
Hence the closure has domain
\[
D(\Hxp)=U^{-1}H^1(\R),
\]
is self-adjoint, and satisfies
\[
\sigma(\Hxp)=\R,\qquad \sigma_p(\Hxp)=\varnothing,
\]
with purely absolutely continuous spectrum.  This is the standard spectral
theorem for momentum on \(L^2(\R)\).

\section*{Extension Audit}

On the full half-line carrier, \(\Hxp\) has deficiency indices \((0,0)\);
there is no one-parameter family of self-adjoint extensions.

The familiar one-parameter family appears only after replacing the full
logarithmic line by a finite logarithmic interval \(I=[a,b]\).  There the
minimal operator \(-i\,d/dt\) has \(U(1)\)-extensions
\[
D_\theta=\{g\in H^1([a,b]):g(b)=e^{i\theta}g(a)\},
\qquad \theta\in[0,2\pi),
\]
with eigenvalues
\[
\lambda_n(\theta)=\frac{2\pi n+\theta}{b-a},\qquad n\in\mathbb Z.
\]
These spectra are arithmetic progressions.  They are not the irregular zeta
ordinate ledger \(\{\gamma_n\}\).  Therefore the finite-box extension family
does not close the Hilbert--Polya route; it is at most a finite-window or
model-spectrum diagnostic.

\section*{Parent Residual}

Define \(\XiHP\) as the typed spectral-defect residual:
\[
\XiHP=0
\quad\Longleftrightarrow\quad
\text{the declared self-adjoint HP operator has point spectrum }
\{\gamma_n\}
\text{ with the zeta multiplicity ledger and no extra spectral residue.}
\]
For the standard \(\Hxp\), this residual cannot vanish because
\(\sigma_p(\Hxp)=\varnothing\).  For a modified Berry--Keating operator with
cutoffs or nonstandard boundary data, \(\XiHP=0\) is precisely the
Hilbert--Polya assertion.  It is therefore target-equivalent unless the
modification and its spectral ledger are supplied independently of the zero
target.

\section*{Initial HP Cascade}

\begin{center}
\begin{tabular}{p{0.18\linewidth}p{0.28\linewidth}p{0.34\linewidth}p{0.15\linewidth}}
\toprule
Branch & Name & Target & Status\\
\midrule
HP-spectrum & standard \(xp\) spectral ledger
& Compare \(\sigma(\Hxp)\) with \(\{\gamma_n\}\).
& bridge failure\\
HP-extension & self-adjoint extension boundary data
& Test the full half-line and finite log-box extension families.
& unique full extension; finite spectra arithmetic\\
HP-modification & Berry--Keating cutoff or modified operator
& Add phase-space cutoffs or boundary data to seek the zeta spectrum.
& target-equivalent without independent records\\
HP-ledger & determinant/trace formula bridge
& Compare a spectral determinant or trace formula to \(\xi\) without target selection.
& no inherited completed ledger bridge\\
\bottomrule
\end{tabular}
\end{center}

\section*{No-Go Transfer}

\begin{center}
\begin{tabular}{p{0.34\linewidth}p{0.14\linewidth}p{0.42\linewidth}}
\toprule
No-go & Transfer & Reason\\
\midrule
public-shadow non-promotion & yes
& formal spectral analogy or finite plots do not promote to a completed zero ledger\\
finite-window Calkin blindness & yes
& finite-box \(xp\) spectra and finite spectral triples require tail/determinant promotion\\
auxiliary-GRH smuggling & yes
& the zeta ordinates cannot be assumed to be a self-adjoint spectrum\\
incomplete character spectrum support-only & retained
& source-side no-go remains for any spectral/source ladder\\
scalar identity not carrier identity & yes
& \(\xi\), an explicit formula, or a determinant analogy does not identify the \(xp\) carrier\\
ledger-relative Hecke non-comparability & retained
& no Hecke-to-HP spectral ledger bridge is supplied\\
Branch C full-carrier ZI-COV(i) foreclosure & retained
& Burnol/Sonine route-local no-go remains separate\\
de Branges target-equivalence / Conrey--Li survival & retained
& de Branges route-local no-go remains separate\\
HP/BK standard \(xp\) bridge failure and modified target-equivalence & new
& standard \(xp\) has continuous spectrum; exact modified matching is HP/RH content\\
\bottomrule
\end{tabular}
\end{center}

\begin{nogo}[HP/BK carrier no-go]
Under the inherited records, the standard Berry--Keating half-line operator
does not provide a spectral carrier for the zeta zero ledger, and modified
operators that would provide exact spectral matching are target-equivalent to
the Hilbert--Polya/RH assertion unless independently constructed.
\end{nogo}

\begin{proof}
The logarithmic unitary \(Uf(t)=e^{t/2}f(e^t)\) sends the standard
symmetrized \(xp\) operator to \(-i\,d/dt\) on \(L^2(\R)\).  The latter is
self-adjoint on \(H^1(\R)\), has purely absolutely continuous spectrum
\(\R\), and has no eigenvalues.  It therefore cannot have the discrete zero
ordinate ledger as point spectrum.

If one instead cuts the logarithmic line to a finite interval, von Neumann
extension theory gives \(U(1)\)-many self-adjoint extensions.  Their spectra
are arithmetic progressions \((2\pi n+\theta)/L\), not the zeta ordinates.
They are finite-window diagnostics, not a completed zeta carrier.

The remaining Berry--Keating possibility is to modify the operator or impose
phase-space boundary data.  But exact equality of the modified spectrum with
\(\{\gamma_n\}\) is the Hilbert--Polya conjecture in operator form.  Without
an independently specified operator, domain, determinant, and completed
ledger bridge, using that equality would be target-equivalence smuggling.
\end{proof}

\section*{CTMT Check}

The HP/BK pivot does not instantiate CTMT as the primary obstruction.  The
standard carrier fails before a carrier-native matrix-element terminus is
reached: its spectral type is continuous.  Matrix entries of a finite-box or
finite spectral approximation would fall under finite-window blindness unless
they are promoted by a completed determinant/ledger theorem.

\section*{Meta-Pattern Status}

\[
\boxed{\mathrm{meta\_pattern\_status}=\text{meta\_pattern\_confirmed}.}
\]

The supporting instances are:
\begin{enumerate}
\item Hecke H6: ledger-relative bridge non-comparability.
\item de Branges \(H(E_{\mathrm{RH}})\): target-equivalence plus
Conrey--Li survival.
\item Hilbert--Polya / Berry--Keating \(H_{xp}\): standard bridge failure and
modified-operator target-equivalence.
\end{enumerate}
The precise meta-pattern is: classical RH-equivalent carriers do not supply
independent closure under inherited records; they terminate as typed no-go
modes unless an independent completed carrier ledger is supplied.

\section*{Classical Theorems Cited}

\begin{itemize}
\item Pólya 1914 / Hilbert--Polya conjecture: self-adjoint spectral
formulation of RH by the zero ordinates.
\item Spectral theorem for \(-i\,d/dt\) on \(L^2(\R)\): continuous momentum
spectrum.
\item von Neumann self-adjoint extension theorem for \(-i\,d/dt\) on a finite
interval: \(U(1)\) boundary family and arithmetic eigenvalues.
\item Berry 1986 and Berry--Keating 1999: \(H=xp\) and phase-space cutoff
proposal.
\item Connes 1999: adelic/noncommutative trace-formula alternative.
\end{itemize}

\section*{Nonclaim Boundary}

This step does not prove RH, does not assume RH, and does not construct a
Hilbert--Polya operator.  It does not close \(\XiHP\), \(\XiBC\), or
\(\XidB\).  It does not claim that every possible modified Berry--Keating
model fails.  It records that the standard \(xp\) carrier has the wrong
spectral type and that exact modified spectral matching is target-equivalent
without independent carrier records.

\end{document}
